至此，我们已讨论了大语言模型与智能体的基础：语言与词元化如何工作，大语言模型如何训练和微调，以及先进推理模型如何推理和自我反思。我们考察了语言、视觉和动作等不同学习模态，也讨论了元学习的复杂之处。

现在，我们来到智能体大语言模型的核心，回顾大语言模型智能体的不同交互方式。本部分包含一章理论内容和一章实践内容。我们先回顾前文涉及的五种智能体交互方式，并讨论记忆的重要性；随后通过一个动手实践案例，把本部分的多项技术结合起来。你将尝试创建一个能够自我反思的大语言模型智能体，让它游玩 TextWorld、NetHack 等典型人工智能游戏。
